\section{Results}\label{sec:results}
% Page budget (12-page CFP limit incl. references): this section <= ~3.3 pages.
% Every number must come from scripts/make_assets.py reading results/.
% Family A and Family B in separate tables/panels; objective attributions from Family B only.

\TODO{Not yet reported. Fill in once Phase 6 full scale (9 models $\times$ 5
logs), Phase 7 (reliability) and Phase 8 (future organisation) are complete.}

\subsection{Sanity and predictive performance}\label{sec:res-sanity}
\TODO{About 0.3 pages. One compact table of validation/test task metrics with
collapse flags (generated). One sentence on the validation-to-test gap under
the chronological split.}

\subsection{RQ1: What do the models learn?}\label{sec:res-rq1}
\TODO{About 1 page. Per model and log: used dimensionality ($\erank/d$, PR),
trajectory shape ($S$, $\theta$, normalised step size, with bootstrap CIs
over cases) and progress profiles (one figure). Describe the patterns that
recur across logs and those that vary with architecture or objective. Make
any objective attribution only from the B1 vs.\ B2 contrast, using the
length-matched comparison. No hypothesis verdict; this question is
descriptive.}

\subsection{RQ2: Does geometry predict errors?}\label{sec:res-rq2}
\TODO{About 1 page. Table of A vs.\ B vs.\ C (AUROC, AUPRC, Brier, $\Delta$
log-likelihood, CIs) per model and log, plus the high-confidence-error
analysis. Give a verdict on H1 with effect sizes, not only $p$-values.}

\subsection{RQ3: Organisation by future}\label{sec:res-rq3}
\TODO{About 1 page. Global alignment, precision@10, trustworthiness, both
history-controlled tests against the non-learned history floor, and one
figure of branch-separation curves over progress. Give verdicts on H2 and
H3.}
